% ============================================================
%  sections/problem.tex
% ============================================================

\section{The Problem}

\subsection{What teams experience today}

Every data engineering team eventually hits this wall:

\begin{enumerate}
  \item A Dagster or Airflow pipeline has been running fine for weeks.
  \item Job latency has been quietly climbing. Nobody notices — there is no alert for "slow trend."
  \item Concurrency increases as the business grows. Resource pressure accumulates.
  \item At 2am, the pipeline fails. Dashboards show stale data. Customers call.
  \item The on-call engineer spends 3–4 hours diagnosing something that was visible in the
        data for two weeks.
\end{enumerate}

\begin{callout}
\textbf{The root cause is not the failure itself — it is the absence of a system that connects
performance trends to infrastructure limits.} Every failure has a predictable precursor.
No current tool surfaces it.
\end{callout}

\subsection{Why existing tools do not solve this}

\begin{center}
\small
\begin{tabularx}{0.96\textwidth}{l X X}
  \toprule
  \textbf{Tool} & \textbf{What it does} & \textbf{What it misses} \\
  \midrule
  Monte Carlo, Bigeye     & Detects data quality anomalies after the fact    & No performance trend; no infra recommendation \\
  Datadog, Grafana        & Infrastructure metrics and alerting              & No pipeline-level semantics; alert fatigue \\
  Great Expectations, dbt & Schema and value assertion testing               & Requires manual rule-writing; misses unknown failures \\
  PagerDuty               & Incident routing after failure                   & Reactive only; no prediction \\
  \textbf{Runway}         & \textbf{Predicts failure date, recommends action} & \textbf{—} \\
  \bottomrule
\end{tabularx}
\end{center}

\vspace{4pt}

\subsection{The gap in one sentence}

\begin{insight}
No tool connects pipeline performance degradation trends to a specific failure date and
a concrete infrastructure recommendation. Runway fills exactly this gap.
\end{insight}
